\chapter{Il tessuto nervoso}

% ---------------------------------------------------------------------------
\section{Generalità}

Il \textbf{tessuto nervoso} è costituito da due grandi popolazioni cellulari:
i \textbf{neuroni}, circa 100 miliardi, e le \textbf{cellule della glia},
che sono ancora più numerose, dell'ordine dei mille miliardi. Il
\textbf{sistema nervoso centrale (SNC)} comprende il cervello e il midollo
spinale; il \textbf{sistema nervoso periferico (SNP)} comprende invece i
nervi e i gangli.

% ---------------------------------------------------------------------------
\section{Il neurone}

I \textbf{neuroni} sono cellule eccitabili, capaci cioè di rispondere a
stimoli fisici o chimici mediante modificazioni della concentrazione di ioni
sulle due facce della membrana plasmatica. In seguito all'eccitazione, il
segnale si manifesta come una variazione del \textbf{potenziale d'azione
(p.a.)} — fenomeno noto come trasduzione — che viene propagato o condotto
lungo la membrana dal punto di insorgenza ad altre parti della cellula. I
neuroni trasmettono i segnali ad altre cellule mediante le \textbf{sinapsi}
e sono inoltre dotate di \textbf{memoria}.

Come ogni altra cellula, il neurone possiede una membrana, un nucleo
contenente i geni, mitocondri e altri organelli, ed è in grado di
sintetizzare proteine e produrre energia. Ciò che lo distingue dalle altre
cellule è la presenza di prolungamenti chiamati \textbf{dendriti} e
\textbf{assoni}: i dendriti portano informazioni verso il corpo cellulare,
mentre gli assoni le portano lontano dal corpo cellulare. I neuroni
contengono inoltre \textbf{neurotrasmettitori} e formano sinapsi.

\subsection{Le quattro regioni funzionali del neurone}

Il neurone può essere suddiviso in quattro regioni funzionali:

\begin{itemize}
  \item i \textbf{dendriti} raccolgono i segnali provenienti dall'ambiente o
    dall'attività di altre cellule; sulla loro superficie sono presenti le
    \textbf{spine dendritiche};
  \item il \textbf{corpo cellulare} contiene il nucleo (a sua volta dotato di
    nucleolo), i mitocondri, i ribosomi e gli altri organelli cellulari; nel
    citoplasma si osservano inoltre i \textbf{corpi di Nissl} e i
    \textbf{neurofilamenti};
  \item l'\textbf{assone}, che origina dal \textbf{monticolo assonico} e
    presenta un \textbf{segmento iniziale}, conduce l'impulso nervoso
    (potenziale d'azione) verso le terminazioni sinaptiche; l'assone può
    essere mielinizzato;
  \item i \textbf{bottoni sinaptici} trasmettono l'impulso ad altri neuroni
    o a effettori quali muscoli e ghiandole.
\end{itemize}

Oltre a nucleo, nucleolo, mitocondri, apparato di Golgi e reticolo
endoplasmatico liscio e rugoso, nel citoplasma del corpo cellulare sono
presenti poliribosomi e ribosomi liberi, microtubuli e microfilamenti che si
estendono anche nell'assone. Il bottone sinaptico contiene vescicole
assonali cariche di neurotrasmettitore, che viene rilasciato nello spazio
della sinapsi e captato da appositi recettori sulla membrana della cellula
postsinaptica, superando la cosiddetta distanza sinaptica. Lungo l'assone,
la guaina mielinica (formata dalle cellule di Schwann) è interrotta a
intervalli regolari dai \textbf{nodi di Ranvier}, e l'assone può inoltre
dare origine a sinapsi assone-assone con il terminale assonale di un altro
neurone.

% ---------------------------------------------------------------------------
\section{Classificazione strutturale dei neuroni}

In base alla morfologia si distinguono quattro tipi di neurone: il neurone
\textbf{anassonico}, privo di un assone identificabile e provvisto di soli
dendriti; il neurone \textbf{bipolare}, che possiede un dendrite e un
assone; il neurone \textbf{pseudounipolare}, il cui corpo cellulare è
provvisto di un unico prolungamento che, a breve distanza dal corpo, si
suddivide a T in un segmento iniziale e in un assone (per questo motivo le
cellule pseudounipolari sono dette anche neuroni bipolari o neuroni a T); il
neurone \textbf{multipolare}, provvisto di un numero vario di dendriti e di
un unico assone. Tutti i neuroni, ad eccezione di quelli anassonici,
possiedono un solo assone: i neuroni unipolari sono provvisti del solo
assone e mancano di dendriti, i neuroni bipolari possiedono un dendrite e un
assone, i neuroni multipolari sono provvisti di un numero vario di dendriti
e dell'assone.

\subsection{Osservazioni istologiche}

Nel midollo allungato del tronco encefalico umano si osservano numerosi
\textbf{neuroni multipolari}: nel citoplasma dei pirenofori è ben visibile
la zona tigroide del Nissl, mentre nei nuclei, quando il piano di taglio
passa esattamente per il centro, si notano chiaramente, con qualunque
colorazione, i nucleoli; i nuclei più piccoli e rotondi appartengono invece
alle cellule della neuroglia (colorazione di Nissl, ingrandimento 200x).

Nella corteccia telencefalica si osserva analogamente un neurone
multipolare, ben evidenziato dalla colorazione di Cajal (ingrandimento
200x).

% ---------------------------------------------------------------------------
\section{Classificazione funzionale dei neuroni}

Dal punto di vista funzionale, il sistema nervoso è organizzato in
recettori, sistema nervoso periferico e sistema nervoso centrale. I
recettori si distinguono in \textbf{interocettori}, \textbf{esterocettori}
e \textbf{propriocettori}, che inviano informazioni tramite \textbf{fibre
afferenti} ai \textbf{neuroni sensitivi} localizzati nei gangli periferici;
questi ultimi trasmettono i segnali al sistema nervoso centrale.

All'interno del SNC, gli \textbf{interneuroni} hanno il compito di
coordinare gli impulsi motori. Il sistema motorio si articola in una
componente \textbf{somatica}, che dai \textbf{neuroni motori somatici}
raggiunge, tramite \textbf{fibre efferenti}, la \textbf{muscolatura
scheletrica}, e in una componente \textbf{viscerale}, che dai
\textbf{neuroni motori viscerali nel SNC} raggiunge, tramite \textbf{fibre
pregangliari}, i \textbf{neuroni motori viscerali nel ganglio periferico
motorio}, i quali proiettano infine, tramite \textbf{fibre postgangliari},
agli \textbf{effettori viscerali}: muscolatura liscia, ghiandole,
muscolatura cardiaca e cellule adipose. Nello schema funzionale, le frecce
di colore blu rappresentano la componente somatica (sensoriale e motoria),
mentre le frecce di colore rosso rappresentano la componente viscerale
(sensoriale e motoria).

% ---------------------------------------------------------------------------
\section{La sinapsi}

La \textbf{sinapsi} è la struttura che permette ai neuroni di connettersi
tra loro: la trasmissione dei segnali elettrici prodotti a livello del
corpo del neurone e propagati lungo l'assone avviene attraverso giunzioni
specifiche costituite fra due neuroni.

In base al bersaglio sinaptico si distinguono la \textbf{sinapsi
asso-somatica}, che si stabilisce tra l'assone di un neurone e il corpo
cellulare di un altro; la \textbf{sinapsi asso-dendritica}, tra l'assone di
un neurone e un dendrite di un altro; e la \textbf{sinapsi asso-assonica},
tra l'assone di un neurone e l'assone di un altro.

\subsection{Sinapsi elettriche e sinapsi chimiche}

Esistono due tipi di sinapsi: le \textbf{sinapsi elettriche}, dette anche
\textit{gap junctions}, in cui il citoplasma del neurone presinaptico e
quello del neurone postsinaptico comunicano direttamente attraverso canali
di giunzione (\textit{gap junction channels}) che attraversano le membrane
presinaptica e postsinaptica, permettendo il passaggio di ioni; questo tipo
di sinapsi può condurre gli impulsi nervosi in entrambe le direzioni. Le
\textbf{sinapsi chimiche}, invece, presentano un neurone presinaptico dotato
di vescicole sinaptiche (\textit{synaptic vesicle}) che, fondendosi con la
membrana presinaptica, rilasciano il neurotrasmettitore nello spazio
sinaptico (\textit{synaptic cleft}); il neurotrasmettitore si lega quindi a
recettori specifici sulla membrana postsinaptica, attivando canali ionici
postsinaptici. Le sinapsi chimiche possono condurre gli impulsi nervosi in
una sola direzione.

Nella sinapsi chimica, tra le due membrane rimane uno spazio in cui sono
localizzate proteine che intervengono nella regolazione del mantenimento
della giunzione e del metabolismo del neurotrasmettitore liberato. Come già
indicato, la sinapsi chimica può aver luogo tra la terminazione assonica e
il dendrite del secondo neurone (sinapsi asso-dendritica), oppure il corpo
del secondo neurone (sinapsi asso-somatica), o infine l'assone del secondo
neurone (sinapsi asso-assonica). Le vescicole sinaptiche sono sferiche e
contengono ciascuna circa 10.000 molecole di neurotrasmettitore.

\subsection{Sinapsi vescicolari}

Gli effetti sulla membrana postsinaptica sono di breve durata, poiché le
molecole di neurotrasmettitore vengono rimosse per azione enzimatica oppure
riassorbite. A livello del corpo di un neurone si osservano numerosi
\textbf{bottoni sinaptici}, connessi tramite l'\textbf{arborizzazione
terminale} agli \textbf{assoni} provenienti da altri neuroni, i cui
prolungamenti sono rivestiti dal \textbf{rivestimento mielinico}; sulla
superficie del neurone sono inoltre visibili i \textbf{dendriti} (spesso
osservati in sezione, cioè "sezionati") e i \textbf{processi gliali} che si
interpongono tra le sinapsi.

% ---------------------------------------------------------------------------
\section{Circuiti neuronali}

I neuroni si organizzano in diversi tipi di circuiti, a seconda della
modalità con cui elaborano e trasmettono l'informazione:

\begin{itemize}
  \item il \textbf{circuito divergente}, in cui un singolo neurone si
    ramifica progressivamente andando a stimolare un numero sempre maggiore
    di neuroni a valle: è il meccanismo alla base della stimolazione diffusa
    di neuroni multipli o di interi pool neuronali nel SNC;
  \item il \textbf{circuito convergente}, in cui più neuroni, con sorgenti
    diverse, convogliano i propri impulsi su un singolo neurone;
  \item il \textbf{circuito di elaborazione seriale}, in cui neuroni singoli
    o gruppi di neuroni lavorano in maniera sequenziale, l'uno dopo l'altro;
  \item il \textbf{circuito di elaborazione in parallelo}, in cui singoli
    neuroni o gruppi di neuroni processano simultaneamente le stesse
    informazioni lungo più vie parallele;
  \item il \textbf{circuito riverberante}, basato su un meccanismo di
    feedback che può essere di tipo eccitatorio o inibitorio, in cui
    l'impulso ritorna, tramite un collaterale, sul neurone di partenza
    prima di proseguire verso il neurone successivo.
\end{itemize}

% ---------------------------------------------------------------------------
\section{Le cellule della glia (nevroglia)}

Le cellule della \textbf{nevroglia} si trovano sia nel sistema nervoso
periferico sia nel sistema nervoso centrale.

Nel \textbf{sistema nervoso periferico} si distinguono:

\begin{itemize}
  \item le \textbf{cellule satelliti}, che circondano i corpi cellulari
    neuronali nei gangli e regolano i livelli di CO\textsubscript{2},
    O\textsubscript{2}, sostanze nutritizie e neurotrasmettitori intorno ai
    neuroni all'interno dei gangli;
  \item le \textbf{cellule di Schwann}, che circondano tutti gli assoni nel
    SNP, sono responsabili della mielinizzazione degli assoni periferici e
    partecipano ai processi di riparazione dopo una lesione.
\end{itemize}

Nel \textbf{sistema nervoso centrale} si distinguono:

\begin{itemize}
  \item gli \textbf{oligodendrociti}, che mielinizzano gli assoni del SNC e
    offrono un supporto strutturale;
  \item gli \textbf{astrociti}, che costituiscono la barriera
    ematoencefalica, offrono un supporto strutturale, regolano la
    concentrazione di ioni, nutrienti e gas, e assorbono e riciclano i
    neurotrasmettitori, formando cicatrici in seguito a lesioni;
  \item la \textbf{microglia}, che rimuove detriti cellulari, rifiuti e
    patogeni per fagocitosi;
  \item le \textbf{cellule ependimali}, che rivestono i ventricoli
    cerebrali e il canale centrale midollare, e regolano la produzione, la
    circolazione e il riassorbimento del liquido cerebrospinale.
\end{itemize}

Nel dettaglio, le \textbf{cellule ependimali} sono organizzate in forma di
epitelio monostratificato per rivestire le cavità interne del SNC
(ventricoli, acquedotto del mesencefalo e canale centrale del midollo
spinale), dove è presente il liquido cerebrospinale (LCR); sono provviste
di microvilli e ciglia unite fra loro da giunzioni.

La \textbf{microglia} viene descritta come una sorta di "colla" del sistema
nervoso: interviene nella regolazione della comunicazione sinaptica e nella
crescita dei neuroni, oltre a intervenire nei processi immunitari.

Gli \textbf{astrociti} sono le cellule della glia più numerose. Hanno forma
irregolarmente stellata, con espansioni citoplasmatiche laminari che si
dispongono a ricoprire la membrana dei corpi neuronali, dei dendriti e delle
giunzioni sinaptiche nella sostanza grigia, trovandosi anche tra i fasci di
fibre mielinizzate della sostanza bianca.

% ---------------------------------------------------------------------------
\section{La mielina}

La \textbf{mielina} migliora la conduzione dell'impulso nervoso. Il
processo di formazione della guaina mielinica, o mielinizzazione, inizia
nel periodo embrionale e si completa in tempi diversi nelle differenti
parti del sistema nervoso; la mielina ha una struttura multilamellare, di
colore biancastro.

Nelle fibre mielinizzate l'impulso nervoso viaggia a velocità sempre
maggiore rispetto alle fibre non mielinizzate, che sono sempre di calibro
inferiore, anche grazie al particolare comportamento "saltatorio"
dell'impulso, che salta da un nodo all'altro. Gli \textbf{oligodendrociti}
possono mielinizzare fino a 60 assoni differenti, mentre le \textbf{cellule
di Schwann} mielinizzano un solo assone.

La guaina mielinica ha la composizione in lipidi e proteine della membrana
plasmatica degli oligodendrociti (o delle cellule di Schwann) e ha un
colore bianco che identifica le parti del sistema nervoso centrale formate
da queste fibre come \textbf{sostanza bianca}.

Nel sistema nervoso periferico il processo di mielinizzazione presenta due
aspetti peculiari: ciascuna cellula di Schwann produce un segmento
internodale di un solo assone; in corrispondenza del nodo di interruzione
della mielina, i lembi citoplasmatici di due cellule di Schwann successive
ricoprono completamente l'assone, isolandolo dal tessuto circostante.

Nella mielina il 40\% è costituito da acqua, il restante da parte solida;
di questa parte solida, l'80\% è costituito da lipidi (colesterolo e
fosfolipidi), a cui si aggiungono proteine e carboidrati. Le cellule di
Schwann si avvolgono intorno al neurone e lo isolano elettricamente.

Uno schema anatomico della fibra mielinica evidenzia, procedendo dall'esterno
verso l'interno: le fibre collagene dell'endonevrio, la lamina basale, la
guaina mielinica, il nucleo di una cellula di Schwann, il nodo di Ranvier e
l'assone.

\subsection{Velocità di conduzione}

Grazie alla mielinizzazione, la propagazione del segnale elettrico avviene
molto più rapidamente, fino a 150 m/s, rispetto agli assoni privi di guaina
mielinica: se l'impulso dovesse infatti percorrere l'intero assone senza
l'ausilio della conduzione saltatoria, la sua velocità si ridurrebbe a soli
5 m/s.

\subsection{Il processo di mielinizzazione}

Durante lo sviluppo di un manicotto mielinico, un assone ancora "nudo"
affonda nel citoplasma di una cellula di Schwann o di un oligodendrocita.
Uno dei due lembi della doccia formatasi nella cellula di Schwann compie
quindi una serie di rotazioni, avvolgendosi a spirale attorno all'assone e
producendo così molteplici sovrapposizioni di strati di membrana
plasmatica. Ad ogni spira della cellula di Schwann o dell'oligodendrocita,
il citoplasma viene rimosso. Al termine di questo processo, il rivestimento
mielinico risulta costituito da un gran numero di doppi strati
fosfolipidici sovrapposti, che svolgono una funzione isolante sulla
membrana dell'assone. Nel sistema nervoso periferico, ciascuna cellula di
Schwann può mielinizzare un solo assone oppure racchiudere, senza
mielinizzarli, numerosi assoni amielinici; nel sistema nervoso centrale, un
singolo oligodendrocita proietta invece prolungamenti che avvolgono e
mielinizzano più assoni contemporaneamente.

% ---------------------------------------------------------------------------
\section{L'assone e la velocità di conduzione}

L'\textbf{assone} è un conduttore di impulsi in direzione centrifuga
rispetto al corpo cellulare. La velocità di conduzione aumenta con
l'aumentare del diametro dell'assone: gli assoni di grande diametro hanno
una minore resistenza ai flussi di corrente longitudinali e quindi una
minore resistenza assiale; di conseguenza essi avranno una maggiore
costante di spazio. In generale, maggiore è il diametro dell'assone,
maggiore è la velocità di propagazione del potenziale d'azione. È quanto si
verifica nelle fibre amieliniche, cioè nelle fibre nervose in cui l'assone
non è provvisto di alcuna forma specializzata di rivestimento.

A titolo di esempio, gli assoni primari afferenti si distinguono in
quattro tipi in base al diametro e alla velocità di conduzione: le fibre
di tipo A\(\alpha\), con diametro di 13-20 \(\mu\)m, conducono a una
velocità di 80-120 m/s; le fibre di tipo A\(\beta\), con diametro di 6-12
\(\mu\)m, conducono a 35-75 m/s; le fibre di tipo A\(\delta\), con diametro
di 1-5 \(\mu\)m, conducono a 5-35 m/s; le fibre di tipo C, prive di
mielina e con diametro di soli 0,2-1,5 \(\mu\)m, conducono a velocità
molto più basse, comprese tra 0,5 e 2,0 m/s.

\subsection{Riassumendo}

Le fibre nervose possono quindi essere distinte in due categorie: le fibre
\textbf{amieliniche}, in cui gli assoni sono semplicemente contenuti in una
nicchia delle cellule di Schwann, senza un vero e proprio rivestimento
mielinico; e le fibre \textbf{mieliniche}, in cui le cellule mieliniche si
avvolgono più volte attorno all'assone, formando una fila di manicotti che
schermano la fibra dalla corrente ionica. Lungo le fibre mieliniche, la
guaina mielinica presenta delle interruzioni, i \textbf{nodi di Ranvier},
che permettono una conduzione del segnale elettrico più rapida, detta
\textbf{conduzione saltatoria}.

% ---------------------------------------------------------------------------
\section{L'assone gigante di calamaro}

Un organismo modello storicamente importante per lo studio della fisiologia
dell'eccitabilità è il \textbf{calamaro}: il suo assone gigante, molto
grande (fino a 1 mm di diametro), controlla parte del sistema di propulsione
a getto d'acqua dell'animale. Scoperto nel 1909, l'assone gigante di
calamaro ha permesso di registrare velocità di conduzione fino a 210 m/s,
tra le più elevate mai misurate.

% ---------------------------------------------------------------------------
\section{Cenni storici: elettricità e stimolazione del sistema nervoso}

Il tessuto nervoso è un \textbf{tessuto eccitabile}, e l'utilizzo di
metodiche di stimolazione elettrica e magnetica per modulare l'attività del
sistema nervoso centrale è noto fin dall'antichità: già nel periodo 43-48
a.C. gli antichi greci e romani furono tra i primi a sfruttare i fenomeni
bioelettrici prodotti da alcuni pesci per la cura di determinate affezioni
neurologiche, come cefalea ed epilessia.

I \textbf{pesci elettrici} — tra cui il pesce gatto elettrico
(\textit{Electric Catfish}), la torpedine (\textit{Electric Ray}),
l'anguilla elettrica (\textit{Electric Eel}), il pesce del Nilo (\textit{Nile
Fish}) e la razza elettrica (\textit{Electric Skate}) — sono in grado di
generare autonomamente campi elettrici tramite appositi organi di origine
muscolare, con funzione di meccanismo di difesa verso i predatori. Un
lavoro storico del 1972 (\textit{Biochemical Journal}, vol. 128, pp.
833-846) ha documentato in dettaglio localizzazione, innervazione e
struttura degli organi elettrici della torpedine, illustrando la posizione
degli occhi (\textit{eyes}), degli spiracoli (\textit{spiracles}) e degli
organi elettrici (\textit{electric organs}) rispetto a encefalo anteriore
(\textit{forebrain}), tetto ottico (\textit{optic tectum}), cervelletto
(\textit{cerebellum}), lobo elettrico della linea laterale (\textit{lateral
line/electric lobe}), nervi diretti all'organo (\textit{nerves to organ}),
placche (\textit{plaques}) e terminazioni nervose (\textit{nerve
terminals}).

Un articolo del gruppo di studio di storia della neurologia (Maggioni,
Mainardi, Dainese, Campagnaro, Zanchin — \textit{Neurological Sciences},
2005, vol. 26, pp. S431-S433) ha inoltre ricostruito le terapie per la
cefalea descritte da Scribonio Largo.

Un riferimento iconografico storico è rappresentato dall'opera di C.
Bischoff, \textit{Commentario de Usu Galvanismi} (1801), che illustra
apparecchiature per l'applicazione del galvanismo alla testa del paziente
tramite elettrodi e bende. Sono inoltre documentate, con fotografie
d'epoca, sale e apparecchiature per l'elettroshock (\textit{electroshock}),
insieme all'immagine, divenuta iconica, dell'applicazione di elettrodi
temporali a un paziente durante una seduta di terapia elettroconvulsivante.

Una rassegna di Hoy e Fitzgerald (\textit{Nature Reviews Neurology}, 2010,
doi:10.1038/nrneurol.2010.30) sulla stimolazione cerebrale in psichiatria e
sui suoi effetti sulla cognizione ripercorre l'evoluzione storica delle
tecniche di neurostimolazione:

\begin{itemize}
  \item \textbf{ECT} (\textit{Electroconvulsive Therapy}, anni '30): induce
    una crisi convulsiva mediante somministrazione di una corrente elettrica
    al cervello, tramite elettrodi posizionati sullo scalpo; il meccanismo
    d'azione esatto non è noto, ma le teorie proposte includono la
    correzione della carenza di neurotrasmettitori inibitori e l'induzione
    di neurogenesi, in particolare a livello dell'ippocampo;
  \item \textbf{tDCS} (\textit{transcranial Direct Current Stimulation},
    anni '60): applica una debole corrente elettrica allo scalpo tramite due
    elettrodi di superficie (anodo e catodo); la tDCS altera l'eccitabilità
    neuronale spostando il potenziale di membrana dei neuroni superficiali
    in direzione depolarizzante o iperpolarizzante; l'effetto dura in
    genere fino a un'ora dopo un singolo periodo di stimolazione;
  \item \textbf{TMS} (\textit{Transcranial Magnetic Stimulation}, anni '90):
    utilizza un campo magnetico focale, intenso (circa 2 T) e variabile nel
    tempo, per indurre un campo elettrico nelle regioni superficiali della
    corteccia; l'induzione del campo magnetico provoca la depolarizzazione o
    la scarica delle cellule nervose nel cervello; se vengono applicati
    treni ripetuti di impulsi, la scarica ripetuta dei neuroni nel tempo
    sembra modificarne l'attività;
  \item \textbf{DBS} (\textit{Deep Brain Stimulation}, 2005): consiste
    nell'impianto diretto di elettrodi in regioni cerebrali localizzate, con
    l'obiettivo di alterare l'attività cerebrale sia locale sia connessa,
    generalmente tramite stimolazione ad alta frequenza; gli elettrodi sono
    collegati a un generatore di impulsi (pacemaker) impiantato sotto la
    cute del torace;
  \item \textbf{EpCS} (\textit{Epidural Cortical Stimulation}, 2007): è una
    tecnica di stimolazione corticale diretta che utilizza elettrodi
    impiantati sopra la dura madre, in corrispondenza della regione
    cerebrale desiderata; i neuroni sottostanti vengono attivati tramite
    induzione di un campo elettrico;
  \item \textbf{100 Hz MST} (\textit{Magnetic Seizure Therapy}, 2008):
    induce una crisi focale tramite TMS ripetitiva ad alta frequenza; la
    crisi ha origine nelle regioni superficiali della corteccia e, a
    differenza dell'ECT, nessuna corrente elettrica attraversa le regioni
    più profonde del cervello; i meccanismi d'azione non sono ancora noti.
\end{itemize}

% ---------------------------------------------------------------------------
\section{Applicazioni sportive della stimolazione cerebrale: la tDCS negli specialisti del salto con gli sci}

Un caso applicativo di particolare interesse per l'ambito sportivo riguarda
l'uso della \textbf{tDCS} (stimolazione transcranica a corrente continua)
negli atleti specialisti del salto con gli sci. Uno studio pubblicato su
\textit{Nature} (17 marzo 2016, vol. 531, p. 283), presentato con il
titolo giornalistico \textit{"Performance boost paves way for 'brain
doping'"}, ha riportato che la stimolazione elettrica sembra migliorare la
resistenza in studi preliminari. Nello specifico, la stimolazione applicata
sulla corteccia motoria in 4 soggetti, confrontata con una stimolazione
placebo (\textit{sham}) applicata in altri 3 soggetti, ha mostrato di:

\begin{itemize}
  \item migliorare le prestazioni;
  \item poter ridurre la capacità di percepire la fatica.
\end{itemize}

Nel dettaglio, la stimolazione ha determinato un miglioramento della forza
di stacco al salto (\textit{jumping force}) del 70\% e della coordinazione
dell'80\%, rispetto al gruppo sottoposto a stimolazione placebo (il
dispositivo di stimolazione utilizzato nello studio era denominato "Halo",
annunciato nel febbraio dello stesso anno).
